\documentclass[a4paper]{article}
\usepackage{iftex}
\ifPDFTeX
  \usepackage{vntex}
\else
  % XeLaTeX / LuaLaTeX (vi du Overleaf chon compiler khac): dung fontspec
  \usepackage{fontspec}
  \setmainfont{TeX Gyre Termes}
  \setsansfont{TeX Gyre Heros}
  \setmonofont{TeX Gyre Cursor}
\fi
%\usepackage[english,vietnam]{babel}
%\usepackage[utf8]{inputenc}

%\usepackage[utf8]{inputenc}
%\usepackage[francais]{babel}
\usepackage{a4wide,amssymb,epsfig,latexsym,multicol,array,hhline,fancyhdr}
\usepackage{booktabs}
\usepackage{longtable}
\usepackage{amsmath}
\usepackage{lastpage}
\usepackage[lined,boxed,commentsnumbered]{algorithm2e}
\usepackage{enumerate}
\usepackage{color}
\usepackage{graphicx}							% Standard graphics package
\usepackage{array}
\usepackage{tabularx, caption}
\usepackage{multirow}
\usepackage[framemethod=tikz]{mdframed}% For highlighting paragraph backgrounds
\usepackage{multicol}
\usepackage{rotating}
\usepackage{graphics}
\usepackage{geometry}
\usepackage{setspace}
\usepackage{epsfig}
\usepackage{tikz}
\usepackage{listings}
\usetikzlibrary{arrows,snakes,backgrounds}
\usepackage{hyperref}
\hypersetup{urlcolor=blue,linkcolor=black,citecolor=black,colorlinks=true} 
%\usepackage{pstcol} 								% PSTricks with the standard color package

\newtheorem{theorem}{{\bf Định lý}}
\newtheorem{property}{{\bf Tính chất}}
\newtheorem{proposition}{{\bf Mệnh đề}}
\newtheorem{corollary}[proposition]{{\bf Hệ quả}}
\newtheorem{lemma}[proposition]{{\bf Bổ đề}}

\everymath{\color{blue}}
%\usepackage{fancyhdr}
\setlength{\headheight}{40pt}
\pagestyle{fancy}
\fancyhead{} % clear all header fields
\fancyhead[L]{
 \begin{tabular}{rl}
    \begin{picture}(25,15)(0,0)
    \put(0,-8){\includegraphics[width=8mm, height=8mm]{logoITSGUsmall.png}}
    %\put(0,-8){\epsfig{width=10mm,figure=hcmut.eps}}
   \end{picture}&
	%\includegraphics[width=8mm, height=8mm]{hcmut.png} & %
	\begin{tabular}{l}
		\textbf{\bf \ttfamily Trường Đại học Sài Gòn}\\
		\textbf{\bf \ttfamily Khoa Công Nghệ Thông Tin}
	\end{tabular} 	
 \end{tabular}
}
\fancyhead[R]{
	\begin{tabular}{l}
		\tiny \bf \\
		\tiny \bf 
	\end{tabular}  }
\fancyfoot{} % clear all footer fields
\fancyfoot[L]{\scriptsize \ttfamily Bài tập lớn 2 (cuối kì) môn Công nghệ phần mềm - Năm học 2026-2027}
\fancyfoot[R]{\scriptsize \ttfamily Trang {\thepage}/\pageref{LastPage}}
\renewcommand{\headrulewidth}{0.3pt}
\renewcommand{\footrulewidth}{0.3pt}


%%%
\setcounter{secnumdepth}{4}
\setcounter{tocdepth}{2}
\makeatletter
\newcounter {subsubsubsection}[subsubsection]
\renewcommand\thesubsubsubsection{\thesubsubsection .\@alph\c@subsubsubsection}
\newcommand\subsubsubsection{\@startsection{subsubsubsection}{4}{\z@}%
                                     {-3.25ex\@plus -1ex \@minus -.2ex}%
                                     {1.5ex \@plus .2ex}%
                                     {\normalfont\normalsize\bfseries}}
\newcommand*\l@subsubsubsection{\@dottedtocline{3}{10.0em}{4.1em}}
\newcommand*{\subsubsubsectionmark}[1]{}
\makeatother

\definecolor{dkgreen}{rgb}{0,0.6,0}
\definecolor{gray}{rgb}{0.5,0.5,0.5}
\definecolor{mauve}{rgb}{0.58,0,0.82}

\lstset{frame=tb,
	language=bash,
	aboveskip=3mm,
	belowskip=3mm,
	showstringspaces=false,
	columns=flexible,
	basicstyle={\small\ttfamily},
	numbers=none,
	numberstyle=\tiny\color{gray},
	keywordstyle=\color{blue},
	commentstyle=\color{dkgreen},
	stringstyle=\color{mauve},
	breaklines=true,
	breakatwhitespace=true,
	tabsize=3,
	numbers=left,
	stepnumber=1,
	numbersep=1pt,    
	firstnumber=1,
	numberfirstline=true
}

\newcommand{\fillme}[1]{{\color{red}[#1]}}
\newcommand{\hinhx}[3]{\begin{center}\includegraphics[width=#2,height=0.66\textheight,keepaspectratio]{#1}\\[2pt]{\small\textit{#3}}\end{center}}
\newcommand{\hinh}[4]{\begin{figure}[htbp]\centering\includegraphics[width=#2,height=0.6\textheight,keepaspectratio]{#1}\caption{#3}\label{#4}\end{figure}}
\newcommand{\hinhngang}[3]{\begin{sidewaysfigure}[p]\centering\includegraphics[width=0.97\textheight,height=0.78\textwidth,keepaspectratio]{#1}\caption{#2}\label{#3}\end{sidewaysfigure}}
\newcommand{\hinhcao}[3]{\begin{figure}[htbp]\centering\includegraphics[height=0.78\textheight,width=\linewidth,keepaspectratio]{#1}\caption{#2}\label{#3}\end{figure}}
\newcommand{\nguon}[1]{\par\vspace{-4pt}{\centering\footnotesize\textit{Nguồn: #1}\par}\medskip}
\usepackage{xurl}
\usepackage{pdflscape}
\usepackage{float}
\renewcommand{\arraystretch}{1.25}
\setlength{\LTpre}{6pt}\setlength{\LTpost}{6pt}
\setlength{\tabcolsep}{4pt}
\newcolumntype{L}[1]{>{\raggedright\arraybackslash}p{\dimexpr #1*93/100\relax}}
\usepackage[htt]{hyphenat}
\newcommand{\code}[1]{\texttt{#1}}
\newcommand{\anh}[3]{\begin{figure}[H]\centering\includegraphics[width=#2,height=0.25\textheight,keepaspectratio]{#1}\caption{#3}\end{figure}}
\newcommand{\haianh}[5]{\begin{figure}[H]\centering\begin{minipage}[t]{0.49\linewidth}\centering\includegraphics[width=\linewidth,height=0.23\textheight,keepaspectratio]{#1}\\[1pt]{\footnotesize #3}\end{minipage}\hfill\begin{minipage}[t]{0.49\linewidth}\centering\includegraphics[width=\linewidth,height=0.23\textheight,keepaspectratio]{#2}\\[1pt]{\footnotesize #4}\end{minipage}\caption{#5}\end{figure}}
\newcommand{\anhdt}[3]{\begin{figure}[H]\centering\includegraphics[height=0.36\textheight,keepaspectratio]{#1}\caption{#3}\end{figure}}
\newcommand{\haianhdt}[4]{\begin{figure}[H]\centering\begin{minipage}{0.45\linewidth}\centering\includegraphics[height=0.38\textheight,keepaspectratio]{#1}\end{minipage}\hfill\begin{minipage}{0.45\linewidth}\centering\includegraphics[height=0.38\textheight,keepaspectratio]{#2}\end{minipage}\caption{#3}\label{#4}\end{figure}}
\captionsetup{font=small}
\graphicspath{{figures/}{uml/}}
\everymath{}
\usepackage{fancyvrb}
\fvset{fontsize=\small,frame=lines,framesep=2mm,numbers=none}
\usepackage{enumitem}
\setlist{itemsep=1pt,parsep=0pt,topsep=3pt}
\setlength{\parskip}{3pt}
\setlength{\emergencystretch}{2em}
\setlength{\intextsep}{8pt plus 2pt minus 2pt}
\setlength{\abovecaptionskip}{4pt}
